\labeledsection{Results}
CONTENT HERE \dots

% \labeledsubsubsection{Potential risks and threads of a plugin-system (wanted dependency injection)}
% Using plugins in an application does come with huge risks.
% One cannot control what an external\footnote{external as in from someone else} \textit{plugin} does or does not do.

% The shown \textit{plugin-system} does not implement any security features and is probably the most basic example of doing dependency injections.
% With the current system, the \textit{plugin} can do virtually everything.
% This includes also accessing the \lstinline{instrumentation} and injecting something else into the classpath, modifying the classpath or exiting/killing the application (\textit{plugin-system}).
% However, with proper rights, a malicious plugin could also access the shell/command-line and do harm to the operating-system.
% While the first scenario is more likely, as \texttt{Java} (/our \textit{plugin-system}) does not implement isolation between \textit{plugins} and the application (/\textit{plugin-system}), the second scenario depends on which rights the application.
% Basically, in which environment the application is running, if it is allowed to access the shell/command-line and if, what it can do there.
% Of course, the \textit{plugin} could also be non-malicious and does exactly what it should do.

% The fact here is that there is no easy way of knowing what a \textit{plugin} might or could do.
% Currently, the second case sounds more dangerous but is actually easier to control: 
% When executing a \textit{plugin-system} it must be run in an \underline{isolated} and has to have as \underline{restricted permissions} as possible.
% Even better if run in a \underline{sandbox} or \underline{container} with next to non non-required features (e.g. removing most of the shell-commands), so that the host-system is fully separated.

% Meanwhile, the first scenario, that the \textit{plugin} modifies the classpath, other classes or returns wrong results, is actually hard to protect against.
% It may be possible to scan \textit{plugins} before loading to ensure that nothing within this JAR is accessing, for example, \texttt{Reflection-Access} / \lstinline{instrumentations} classes.
% However, this would require an entirely new project, including the research needed for this.
% This also would probably restrict the overall functions of \textit{plugins} heavily.

% The only truly safe way is to trust a developer, test \textit{plugins} in a protected environment first and, if available, look into the source code or decompile\footnote{Translating Java-ByteCode back to Java/Human Readable-Code} it.
% Additionally, the entire environment around the application has to be kept clean and secure to ensure that no other person can install an unwanted or malicious \textit{plugin}.
% However, all this requires external solutions.